% ===== PRÉAMBULE CANONIQUE — à coller VERBATIM en tête de CHAQUE .tex =====
\documentclass[11pt,a4paper]{article}
\usepackage[utf8]{inputenc}\usepackage[T1]{fontenc}\usepackage{lmodern}
\usepackage[french]{babel}
\usepackage{amsmath,amssymb,mathtools}\usepackage{icomma}
\usepackage[version=4]{mhchem}          % \ce{...} : équations & formules
\usepackage{siunitx}\sisetup{locale=FR} % \qty{}{}, \si{}, \ang{}
\usepackage{chemfig}                    % \chemfig{...} : structures développées
\usepackage{xcolor}\usepackage{colortbl}
\usepackage{tikz}\usepackage{pgfplots}\pgfplotsset{compat=1.18}
\usepackage[most]{tcolorbox}\usepackage{enumitem}\usepackage{array,booktabs}
\usepackage[a4paper,margin=2.2cm]{geometry}
\usetikzlibrary{arrows.meta,calc,angles,quotes,positioning,patterns,backgrounds,shapes.geometric}
\definecolor{primary}{RGB}{222,49,99}\definecolor{primarydark}{RGB}{150,22,55}
\definecolor{defcol}{RGB}{0,114,178}\definecolor{methcol}{RGB}{52,92,148}
\definecolor{excol}{RGB}{0,135,90}\definecolor{attncol}{RGB}{200,40,55}
\tcbset{boxbase/.style={enhanced,breakable,boxrule=0.9pt,arc=2.5pt,
  left=3mm,right=3mm,top=2.3mm,bottom=2.3mm,fonttitle=\bfseries,coltitle=white}}
% --- boîtes TUTORIEL (noms EXACTS, ne pas renommer) ---
\newtcolorbox{cle}[1][]{boxbase,colback=defcol!6,colframe=defcol,
  title={\textsc{À retenir}\ifx\relax#1\relax\relax\else\textmd{ : \itshape #1}\fi}}
\newtcolorbox{methode}[1][]{boxbase,colback=methcol!7,colframe=methcol,sharp corners,
  title={\textsc{Méthode}\ifx\relax#1\relax\relax\else\textmd{ --- \itshape #1}\fi}}
\newtcolorbox{exemplebox}[1][]{boxbase,colback=excol!5,colframe=excol,
  title={\textsc{Exemple}\ifx\relax#1\relax\relax\else\textmd{ : \itshape #1}\fi}}
\newtcolorbox{piege}[1][]{boxbase,colback=attncol!6,colframe=attncol,
  title={\textsc{Piège}\ifx\relax#1\relax\relax\else\textmd{ : \itshape #1}\fi}}
\newtcolorbox{remarque}[1][]{boxbase,colback=black!4,colframe=black!45,
  title={\textsc{Remarque}\ifx\relax#1\relax\relax\else\textmd{ : \itshape #1}\fi}}
% --- boîtes EXERCICE / CORRIGÉ (noms EXACTS, auto-comptées) ---
\newtcolorbox[auto counter]{exo}[1][]{boxbase,colback=primary!4,colframe=primary,
  title={\textsc{Exercice}~\thetcbcounter\ifx\relax#1\relax\relax\else\textmd{ --- \itshape #1}\fi}}
\newtcolorbox[auto counter]{corrige}[1][]{boxbase,colback=excol!4,colframe=excol,
  title={\textsc{Corrigé}~\thetcbcounter\ifx\relax#1\relax\relax\else\textmd{ --- \itshape #1}\fi}}
\newcommand{\R}{\mathbb{R}}\newcommand{\N}{\mathbb{N}}\newcommand{\Z}{\mathbb{Z}}\newcommand{\ds}{\displaystyle}
% (un \usepackage{fancyhdr}+en-tête est possible ; il est ignoré côté web)
% =========================================================================
\begin{document}
\newcommand{\CF}{\ensuremath{\mathrm{C.F.}}}
\begin{corrige}[appliquer la formule]
$\CF = N_v - 2\times(\text{doublets libres}) - (\text{liaisons})$.
\begin{enumerate}[label=\alph*)]
  \item $\CF = 6 - 2\times 2 - 2 = \mathbf{0}$.
  \item $\CF = 6 - 2\times 3 - 1 = \mathbf{-1}$.
  \item $\CF = 5 - 2\times 1 - 3 = \mathbf{0}$.
  \item $\CF = 5 - 2\times 0 - 4 = \mathbf{+1}$.
  \item $\CF = 4 - 2\times 0 - 4 = \mathbf{0}$.
\end{enumerate}
\end{corrige}

\begin{corrige}[molécules neutres]
Dans une molécule neutre, chaque atome dans sa valence habituelle a une charge formelle nulle.
\begin{enumerate}[label=\alph*)]
  \item \ce{O} : $6 - 4 - 2 = \mathbf{0}$.
  \item \ce{N} : $5 - 2 - 3 = \mathbf{0}$.
  \item \ce{C} : $4 - 0 - 4 = \mathbf{0}$.
  \item \ce{Cl} : $7 - 6 - 1 = \mathbf{0}$.
  \item \ce{H} : $1 - 0 - 1 = \mathbf{0}$.
\end{enumerate}
\end{corrige}

\begin{corrige}[atomes chargés]
\begin{enumerate}[label=\alph*)]
  \item \ce{N} : $5 - 0 - 4 = \mathbf{+1}$ (d'où la charge $+$ de \ce{NH4+}).
  \item \ce{O} : $6 - 2 - 3 = \mathbf{+1}$ (d'où la charge $+$ de \ce{H3O+}).
  \item \ce{O} : $6 - 6 - 1 = \mathbf{-1}$.
  \item \ce{B} : $3 - 0 - 4 = \mathbf{-1}$.
\end{enumerate}
\end{corrige}

\begin{corrige}[la somme redonne la charge]
On additionne toutes les charges formelles.
\begin{enumerate}[label=\alph*)]
  \item \ce{CO3^2-} : $0 + 0 + (-1) + (-1) = \mathbf{-2}$. \checkmark
  \item \ce{NH4+} : $(+1) + 4\times 0 = \mathbf{+1}$. \checkmark
  \item \ce{NO3-} : $(+1) + 0 + (-1) + (-1) = \mathbf{-1}$. \checkmark
\end{enumerate}
\end{corrige}

\begin{corrige}[interpréter le signe]
La \CF{} compare les électrons « propres » de l'atome à ceux de l'atome neutre isolé.
\begin{enumerate}[label=\alph*)]
  \item $\CF = +1$ : l'atome a \textbf{un électron de moins} (il en a « perdu » un dans la structure).
  \item $\CF = -1$ : l'atome a \textbf{un électron de plus}.
  \item $\CF = 0$ : l'atome a \textbf{autant} d'électrons que l'atome neutre (situation idéale).
\end{enumerate}
\end{corrige}

\end{document}
